% Generated draft: requires mathematical and source review.
\documentclass[11pt]{article}
\usepackage[margin=25mm]{geometry}
\usepackage{fontspec}
\setmainfont{Latin Modern Roman}
\usepackage{amsmath,amssymb,mathtools,unicode-math}
\setmathfont{Latin Modern Math}
\usepackage{enumitem,needspace,xurl,hyperref}
\setlength{\parindent}{0pt}
\setlength{\parskip}{5pt}
\setlength{\emergencystretch}{3em}
\newcommand{\N}{\mathbb{N}}
\newcommand{\Z}{\mathbb{Z}}
\newcommand{\Q}{\mathbb{Q}}
\newcommand{\R}{\mathbb{R}}
\newcommand{\C}{\mathbb{C}}
\newcommand{\F}{\mathbb{F}}
\newcommand{\A}{\mathbb{A}}
\newcommand{\PP}{\mathbb{P}}
\newcommand{\E}{\mathbb{E}}
\newcommand{\eps}{\varepsilon}
\newcommand{\dd}{\,\mathrm{d}}
\begin{document}
\section*{UnsolvedMath ID: 30006116}
\textbf{Source problem code:} catalogue entry
\subsection{Definitions and mathematical statement}
The supplied record concerns the intended setting of representations of \(p\)-adic groups. Write \(G(F)\) for the group of \(F\)-points under consideration, \(\operatorname{Irr}_{\mathrm{sm}}(G(F))\) for its irreducible smooth representations, \(W_F\) for the Weil group, and \(G^\vee\) for the dual group. The precise hypotheses on \(F\), \(G\), and representation coefficients are not supplied.

Let \((X,r)\), with \(r>0\), be a Moy--Prasad datum attached to \(\pi\in\operatorname{Irr}_{\mathrm{sm}}(G(F))\) in the sense of the source. The excerpt attributes to Chen--DeBacker--Tsai a construction of a continuous homomorphism
\[
\phi_X:W_F^{=r}\longrightarrow G^\vee,
\]
depending only on the ``associativity class'' of \(X\) and a choice of \(\psi_F\). Neither that class nor the normalization and nature of \(\psi_F\) is defined in the supplied excerpt.

\textbf{Conjectural compatibility.} For every representation and attached positive-depth datum covered by the intended construction, and its corresponding Langlands parameter \(\phi_\pi\), the proposed general compatibility is
\[
\left.\phi_\pi\right|_{W_F^{>r}}=1,
\qquad
\left.\phi_\pi\right|_{W_F^{=r}}=\phi_X.
\]
Here the existence and normalization of the correspondence assigning \(\phi_\pi\) are prerequisites whose general scope is not specified by the record.

The notation in the second equality needs clarification. Under the usual possible interpretation using upper-numbering ramification groups,
\[
W_F^{>r}=W_F^{r+},\qquad
W_F^{=r}=W_F^r/W_F^{r+}.
\]
With this interpretation, the second displayed equality means that, after the first equality makes descent possible, the map induced by \(\left.\phi_\pi\right|_{W_F^r}\) on \(W_F^r/W_F^{r+}\) equals \(\phi_X\). This interpretation is conditional: the supplied excerpt does not define these symbols. It also does not specify the identification of the relevant parameter target with \(G^\vee\), or whether equality means equality of chosen representatives or equality up to conjugacy.

The quoted source asserts this compatibility when \(\phi_\pi\) is a supercuspidal Langlands parameter of Kaletha, referring to \([\mathrm{Kal19},\mathrm{Kal21}]\), and the datum is attached to \(\pi\) ``like before.'' The requested conjecture extends that compatibility beyond the stated case. The omitted preceding hypotheses and attachment procedure are essential unresolved qualifications; this draft does not assert a theorem for arbitrary fields, groups, or choices of datum.
\subsection{Short English statement}
The proposed conjecture says that a representation's Langlands parameter is trivial above the positive depth \(r\), and that its depth-\(r\) part agrees with the parameter constructed from its Moy--Prasad datum. The supplied excerpt states this compatibility for specified Kaletha supercuspidal parameters. A precise general statement still requires the omitted hypotheses, ramification conventions, and parameter normalizations.
\subsection{Sources}
\begin{itemize}
\item Local UnsolvedMath registry, permanent ID \(30006116\), title \emph{Moy-Prasad Data and Langlands Parameters}, exact source code \texttt{catalogue entry}, in the Oberwolfach Reports Open Problems collection. No individual proposer is supplied; the recorded proposed year is 2024.
\item Supplied report citation: \emph{Representations of \(p\)-adic Groups (2025)}, Oberwolfach Reports, pp.~3103--3178, DOI \texttt{10.4171/owr/2024/54}. \url{https://doi.org/10.4171/owr/2024/54}. The report title's year is reproduced as supplied alongside the DOI's 2024 component.
\item Supplied literature source, ``OWR report'': \url{https://publications.mfo.de/bitstream/handle/mfo/4270/OWR_2024_54.pdf?sequence=1}.
\item Supplied literature source, ``Chen--DeBacker--Tsai'': \url{https://arxiv.org/abs/2509.07780}. The excerpt attributes the construction of \(\phi_X\) and the stated compatibility case to Chen--DeBacker--Tsai.
\item The excerpt cites Kaletha as \([\mathrm{Kal19},\mathrm{Kal21}]\); full bibliographic entries are not supplied.
\item Literature assessment dated exactly \texttt{2026-08-22}: the supplied curation labels the problem open and states that the 2025 Moy--Prasad-types work constructs restricted parameters and proves compatibility in specified cases, while general compatibility remains conjectural. This is supplied dated curation, not an independently verified literature assessment.
\item The registry labels the statement \texttt{corrected\_verified} and reports that it was reconstructed and checked against the official Oberwolfach report. These are supplied provenance claims, not independent verification performed for this draft.
\end{itemize}
\end{document}
